\begin{table}[htbp]
\centering
\caption{標記之真值}
\label{tab:groundTrue}
\begin{adjustbox}{max width=\textwidth}
\renewcommand{\arraystretch}{1.4}
\begin{tabular}{>{\centering\arraybackslash}m{1cm} >{\centering\arraybackslash}m{16cm} }
\toprule
項次 & 內容 \\
\toprule
1-1 & ("在", "新北市", None), (None, "瑞芳區", None), (None, "萬瑞快速道路", "上"), ("西行", None, None), (None, "瑞濱端", None), (None, "龍潭堵隧道", "內") \\
\hline
1-2 & (None, None, "北上"), ("在", "中山高速公路", "上"), (None, "10.5公里", "處"), (None, "汐止交流道", None) \\
\hline
1-3 & ("在", "臺北市",  None), (None, "大安區", None), (None, "和平東路二段118巷2弄", None), (None, "科技大樓", "後方"), (None, "科技部", "附近") \\
\hline
1-4 & ("在", "臺北市",  None), (None, "大同區", None), (None, "民權西路", "路口"), (None, "承德路二段", "路口"), (None, "捷運民權西路站", "附近"), (None, "台灣銀行", "附近"), (None, "台北大橋", "前") \\
\hline
1-5 & ("在", "臺北市",  None), (None, "大安區", None), (None, "建國南北快速道路", None), (None, "忠孝東路三段", "上方"), ("鄰近", "臺北科技大學", None) \\
\hline
1-6 & (None, "台北大橋",  None), ("往", "台北", None), ("在", "新店溪", "上") \\
\hline
2-1 & (None, "新店溪", None), (None, "新店溪", "沿岸"), (None, "翡翠水庫", "下游")\\
\hline
2-2 & (None, "大漢溪", None), (None, "淡水河", None), (None, "新店溪", "沿岸"), (None, "淡水河", "沿岸"), (None, "石門水庫", "下游")\\
\hline
2-3 &  (None, "木瓜溪", "下游"), (None, "木瓜溪", "沿岸") \\
\hline
3-1 & (None, "台灣", "東北部") \\
\hline
3-2 & (None, "台灣", "東北部"),  (None, "花蓮", "地區") \\
\hline
3-3 & (None, "台灣", "東北部"),  (None, "宜花", "地區") \\
\hline
4-1 & (None, "烏來區", None), (None, "石碇區", None), (None, "南勢溪", "上游"), (None, "崩山溪", "上游")\\
\hline
4-2 & (None, "基隆市", "山區"), (None, "五分山", None)\\
\hline
4-3 & (None, "七堵區", None), (None, "暖暖區", None), (None, "八堵區", None)\\
\hline
5-1 &  (None, "東北部", "沿海"), (None, "宜蘭縣", "沿海")\\
\hline
5-2 &  (None, "北部", "沿海") \\
\hline
5-3 & (None, "東部", "沿海"), (None, "宜花東", "沿海") \\
\bottomrule
\end{tabular}
\end{adjustbox}
\end{table}